% !TeX root = ../机器学习教材（李航）笔记.tex
% 由主文档通过 \input 载入；请编译主文档。
\setcounter{chapter}{8} % 保留原书章号。
\chapter{EM 算法及其推广}
\label{chap:lihang-em}

在混合分布、缺失数据和潜在类别等问题中，观测数据的分布往往需要对未观测变量求和或积分。直接极大化其似然可能很困难，而假如这些变量也已知，参数估计反而容易。\emph{期望极大算法}（expectation maximization algorithm，EM 算法）利用这一差别：先根据当前模型计算隐变量的条件分布，再极大化完全数据对数似然的条件期望。本章对应教材第 9 章，依次讨论算法的构造、单调性与收敛条件、高斯混合模型中的具体更新，以及基于辅助函数的推广。

本章以 $Y$ 表示观测随机变量或观测随机向量，$y$ 表示已经取得并固定的数据；以 $Z$ 表示隐变量，$z$ 表示其可能取值。参数为 $\theta\in\Theta\subseteq\mathbb R^d$，$p_\theta(y,z)$ 统一表示联合概率质量函数或概率密度函数。除特别说明外，求和形式的理论推导先在有限隐状态、关于参数不变的正支撑上进行，并要求观测似然为正、涉及的期望有限；连续隐变量的情形需将求和换成积分，并另行保证可积性及交换微分与积分的条件。

\section{EM 算法的引入}
\label{sec:lihang-em-introduction}

\subsection{EM 算法}
\label{subsec:lihang-em-algorithm}

\emph{观测变量}（observable variable）可由数据取得，\emph{隐变量或潜在变量}（hidden variable or latent variable）则在所用模型中存在而没有被观测。二者共同组成的 $(y,z)$ 称为\emph{完全数据}（complete data），只有 $y$ 时称为\emph{不完全数据}（incomplete data）。这些名称相对于所选模型而言，并不表示观测数据本身存在记录错误。参数估计的目标仍然是观测数据的极大似然估计（maximum likelihood estimation，MLE）：
\begin{equation}
  p_\theta(y)=\sum_z p_\theta(y,z),\qquad
  \ell(\theta)=\log p_\theta(y),\qquad
  \widehat\theta\in\operatorname*{arg\,max}_{\theta\in\Theta}\ell(\theta).
  \label{eq:lihang-em-observed-objective}
\end{equation}
主要计算困难常来自对数内部的求和；这并不意味着一切含隐变量的模型都没有解析解，也不意味着把隐变量直接补成一个确定值就能得到正确估计。

\begin{example}
\label{ex:lihang-em-three-coins}
教材例 9.1 的三硬币模型。硬币 A、B、C 的正面概率分别为 $\pi,p,q$。每次先掷 A，正面选 B，反面选 C，再掷所选硬币，只记录第二次是否为正面。独立重复 $n=10$ 次，观测序列为
\[
  y=(1,1,0,1,0,0,1,0,1,1).
\]
建立概率模型，推导 EM 更新，并说明这些数据能够确定哪些参数信息。
\end{example}

\begin{solve}
令 $Z_i=1$ 表示第 $i$ 次选 B，$Z_i=0$ 表示选 C，$Y_i\in\{0,1\}$ 表示记录结果。模型假设为
\[
  P_\theta(Z_i=1)=\pi,\qquad
  P_\theta(Y_i=1\mid Z_i=1)=p,\qquad
  P_\theta(Y_i=1\mid Z_i=0)=q,
\]
其中 $\theta=(\pi,p,q)\in[0,1]^3$。边缘化 $Z_i$，得到
\begin{equation}
  p_\theta(y_i)=\pi p^{y_i}(1-p)^{1-y_i}
  +(1-\pi)q^{y_i}(1-q)^{1-y_i},\qquad y_i\in\{0,1\}.
  \label{eq:lihang-em-coin-marginal}
\end{equation}
若每次选了哪枚硬币也已知，完全数据的对数似然就是
\[
  \ell_{\mathrm c}(\theta;y,z)
  =\sum_{i=1}^n\bigl\{
  z_i[\log\pi+y_i\log p+(1-y_i)\log(1-p)]
  +(1-z_i)[\log(1-\pi)+y_i\log q+(1-y_i)\log(1-q)]\bigr\}.
\]
它将三个参数分别放入伯努利分布的对数似然中，因而容易极大化。EM 用条件期望替代这里未知的 $z_i$。在第 $t$ 轮参数 $\theta^{(t)}$ 下，根据贝叶斯公式计算
\begin{equation}
  r_i^{(t)}=P_{\theta^{(t)}}(Z_i=1\mid Y_i=y_i)
  =\frac{\pi^{(t)}(p^{(t)})^{y_i}(1-p^{(t)})^{1-y_i}}
  {\pi^{(t)}(p^{(t)})^{y_i}(1-p^{(t)})^{1-y_i}
   +(1-\pi^{(t)})(q^{(t)})^{y_i}(1-q^{(t)})^{1-y_i}}.
  \label{eq:lihang-em-coin-responsibility}
\end{equation}
由于样本对 $(Y_i,Z_i)$ 相互独立，给定全部观测时该条件概率也只依赖 $y_i$。固定这些概率，将 $\ell_{\mathrm c}$ 中的 $z_i$ 换成 $r_i^{(t)}$，便得到待极大化的条件期望。对 $\pi$ 求导的方程为
\[
  \frac{\sum_i r_i^{(t)}}{\pi}
  -\frac{\sum_i(1-r_i^{(t)})}{1-\pi}=0;
\]
对 $p,q$ 分别求导同样得到加权频率。因此新参数为
\begin{equation}
  \begin{aligned}
    \pi^{(t+1)}&=\frac1n\sum_{i=1}^n r_i^{(t)},\\
    p^{(t+1)}&=\frac{\sum_{i=1}^n r_i^{(t)}y_i}{\sum_{i=1}^n r_i^{(t)}},\qquad
    q^{(t+1)}=\frac{\sum_{i=1}^n(1-r_i^{(t)})y_i}{\sum_{i=1}^n(1-r_i^{(t)})}.
  \end{aligned}
  \label{eq:lihang-em-coin-update}
\end{equation}
三个加权伯努利对数似然均为凹函数；相应的成功与失败权重都为正时，上述驻点就是各自唯一的最大点。若某一成分的总权重为零，该成分的正面概率无法由本轮目标确定；若只有成功或只有失败权重，则最大值在参数边界上。取三个参数的初值都严格位于 $(0,1)$ 内，可使本例两种观测下的响应概率均有定义。

从 $(0.5,0.5,0.5)$ 出发，所有 $r_i^{(0)}=0.5$，第一轮得到 $(0.5,0.6,0.6)$，此后保持不变。从 $(0.4,0.6,0.7)$ 出发，第一轮得到约 $(0.406417,0.536842,0.643243)$，也已达到本例的一个固定点。理解这些不同答案，需要注意观测分布实际上只是正面概率为
\[
  \rho=\pi p+(1-\pi)q
\]
的伯努利分布。这十次观测中有六次正面，故
\[
  \ell(\theta)=6\log\rho+4\log(1-\rho),\qquad
  \frac{\mathrm d\ell}{\mathrm d\rho}=\frac6\rho-\frac4{1-\rho}.
\]
该函数在 $0<\rho<1$ 上严格凹，最大点为 $\widehat\rho=0.6$。所以满足 $\pi p+(1-\pi)q=0.6$ 的全部可行三元组都是极大似然解；观测分布不能唯一确定 $\pi,p,q$，称为\emph{参数不可辨识}（non-identifiability）。这里可直接写出极大似然解集，采用 EM 的意义在于展示隐变量迭代机制。特别地，$\widehat\pi=0.5$ 是某次初始化所得解的一部分，不能据此断言数据证明了硬币 A 均匀。
\end{solve}

\begin{definition}[$Q$ 函数]
\label{def:lihang-em-q-function}
固定观测 $y$ 和当前参数 $\theta^{(t)}$，定义
\begin{equation}
  Q(\theta\mid\theta^{(t)})
  =\mathbb E_{\theta^{(t)}}[\log p_\theta(y,Z)\mid Y=y]
  =\sum_z p_{\theta^{(t)}}(z\mid y)\log p_\theta(y,z).
  \label{eq:lihang-em-q-function}
\end{equation}
它是完全数据对数似然关于当前隐变量后验分布的期望。第一个参数是待优化变量，竖线后的参数只确定本轮使用的条件分布。
\end{definition}

一般 EM 算法的输入是观测 $y$、联合模型 $p_\theta(y,z)$、可行参数域 $\Theta$，以及计算隐变量条件分布的方法；隐变量的实际取值不是已知输入。选择使 $p_{\theta^{(0)}}(y)>0$ 且期望存在的初值，规定正容差 $\varepsilon$ 与最大迭代次数 $T$，然后进行如下迭代：
\begin{enumerate}
  \item \emph{E 步}（expectation step）：利用当前后验 $p_{\theta^{(t)}}(z\mid y)$，构造式~\eqref{eq:lihang-em-q-function} 中的 $Q(\theta\mid\theta^{(t)})$。
  \item \emph{M 步}（maximization step）：在可行域内求
  \begin{equation}
    \theta^{(t+1)}\in\operatorname*{arg\,max}_{\theta\in\Theta}
    Q(\theta\mid\theta^{(t)}).
    \label{eq:lihang-em-m-step}
  \end{equation}
  \item 计算观测对数似然。若
  \[
    |\ell(\theta^{(t+1)})-\ell(\theta^{(t)})|
    \leq\varepsilon\bigl(1+|\ell(\theta^{(t)})|\bigr),
  \]
  则按所定精度停止；达到 $T$ 时也终止，并记录是否满足容差。输出最后的参数及其目标值。
\end{enumerate}
式~\eqref{eq:lihang-em-m-step} 的极大值需要存在且能够求得；不能达到精确极大值时，可使用第~\ref{subsec:lihang-em-gem} 小节的广义更新。E 步求的是关于隐变量的条件期望，通常仍是 $\theta$ 的函数，并非只计算一个数；它也不是先把隐变量固定成概率最大的类别。停止条件用于控制数值计算，似然变化很小并不构成全局最优的证明。若采用参数变化准则，还需考虑各参数的尺度。

\begin{notealert}[type=tip,title={从完全数据构造算法}]
推导具体 EM 算法时，先选定隐变量并写出联合模型，再计算完全数据对数似然，最后求其中隐变量相关项的后验期望。若目标是极大后验估计（maximum a posteriori estimation，MAP），对给定先验密度 $p(\theta)$，M 步改为极大化 $Q(\theta\mid\theta^{(t)})+\log p(\theta)$；E 步仍使用给定当前参数的隐变量后验，单调性所针对的目标则是 $\ell(\theta)+\log p(\theta)$。
\end{notealert}

\subsection{EM 算法的导出}
\label{subsec:lihang-em-derivation}

EM 的关键不是让完全数据和观测数据的似然成为同一个函数，而是构造一个在当前参数处与目标相等的下界。记本轮隐变量后验为 $w_t(z)=p_{\theta^{(t)}}(z\mid y)$。在共同正支撑的假设下，用对数函数的凹性，即 Jensen 不等式，得到
\begin{equation}
  \begin{aligned}
    \ell(\theta)
    &=\log\sum_z w_t(z)\frac{p_\theta(y,z)}{w_t(z)}\\
    &\geq\sum_z w_t(z)\log\frac{p_\theta(y,z)}{w_t(z)}
    =:B(\theta\mid\theta^{(t)}).
  \end{aligned}
  \label{eq:lihang-em-jensen-bound}
\end{equation}
这里 $w_t(z)\geq0$ 且 $\sum_z w_t(z)=1$。令
\[
  H(w_t)=-\sum_z w_t(z)\log w_t(z),
\]
则下界可以写成
\begin{equation}
  B(\theta\mid\theta^{(t)})=Q(\theta\mid\theta^{(t)})+H(w_t),\qquad
  B(\theta^{(t)}\mid\theta^{(t)})=\ell(\theta^{(t)}).
  \label{eq:lihang-em-touching-bound}
\end{equation}
第二个等式来自 $p_{\theta^{(t)}}(y,z)/w_t(z)=p_{\theta^{(t)}}(y)$：在当前参数处，Jensen 不等式中的被取平均量不随 $z$ 变化，因而取等号。固定 $w_t$ 后，$H(w_t)$ 与待优化参数 $\theta$ 无关，所以极大化 $B$ 正好等价于 M 步极大化 $Q$。

每一轮都先在当前点建立一个接触目标的下界，再提高这个下界，随后用新参数重新建立下界。E 步的期望若可精确求出，就不涉及随机近似；算法对于原目标的局限在于这种逐轮提高下界的过程通常只找到某个候选解，未必完成原似然的全局极大化。用 $B$ 的几何解释时，必须同时保留下界关系与当前点取等号这两个条件，单独提高一个任意下界不足以推出目标值提高。

\subsection{EM 算法在无监督学习中的应用}
\label{subsec:lihang-em-unsupervised}

无监督学习没有给出每个观测对应的类别或其他目标输出。设输入为 $X$，潜在类别为 $C$，生成模型为 $p_\theta(x,c)$。假设各样本独立同分布，则由观测样本 $x_1,\ldots,x_n$ 估计参数的目标是
\[
  \sum_{i=1}^n\log p_\theta(x_i)
  =\sum_{i=1}^n\log\sum_c p_\theta(x_i,c).
\]
将未给出的类别看作隐变量，就得到 EM 可以处理的形式。E 步给每个可能类别分配后验权重，M 步用这些权重重新估计生成模型；高斯混合模型和朴素贝叶斯模型的无监督学习都是具体例子。EM 是参数估计方法，其适用范围也包括缺失数据等问题；无监督任务能否采用 EM，取决于所设联合模型以及 E、M 两步是否可计算。

\section{EM 算法的收敛性}
\label{sec:lihang-em-convergence}

\subsection{观测似然的单调性}

\begin{theorem}[提高 $Q$ 函数足以保持似然单调]
\label{thm:lihang-em-monotonicity}
在本章下界构造的条件下，若一次参数更新满足
\[
  Q(\theta^{(t+1)}\mid\theta^{(t)})\geq Q(\theta^{(t)}\mid\theta^{(t)}),
\]
则 $\ell(\theta^{(t+1)})\geq\ell(\theta^{(t)})$，从而观测似然也不减。精确的 EM 更新满足上述条件。
\end{theorem}

\begin{proof}
由式~\eqref{eq:lihang-em-jensen-bound} 和式~\eqref{eq:lihang-em-touching-bound}，
\begin{equation}
  \begin{aligned}
    \ell(\theta^{(t+1)})
    &\geq B(\theta^{(t+1)}\mid\theta^{(t)})\\
    &=Q(\theta^{(t+1)}\mid\theta^{(t)})+H(w_t)\\
    &\geq Q(\theta^{(t)}\mid\theta^{(t)})+H(w_t)
    =\ell(\theta^{(t)}).
  \end{aligned}
  \label{eq:lihang-em-monotonicity-chain}
\end{equation}
精确 M 步的可行候选中包含 $\theta^{(t)}$，所以其最大值不会小于当前值。指数函数严格递增，再将对数似然的不等式取指数即可。
\end{proof}

还可以把每次增量分解，以看清何时取等号。对两个有限状态分布 $a,b$，定义\emph{相对熵或 KL 散度}（Kullback--Leibler divergence）
\[
  D_{\mathrm{KL}}(a\Vert b)=\sum_z a(z)\log\frac{a(z)}{b(z)}.
\]
约定 $0\log(0/b)=0$，若 $a(z)>0$ 而 $b(z)=0$，则散度为 $+\infty$。由 $\log u\leq u-1$，在散度有限时有
\[
  -D_{\mathrm{KL}}(a\Vert b)
  =\sum_{z:a(z)>0}a(z)\log\frac{b(z)}{a(z)}
  \leq\sum_{z:a(z)>0}b(z)-1\leq0.
\]
因此散度非负，且当且仅当两个分布相同才为零。定义教材证明所用的条件对数概率期望
\[
  H(\theta\mid\theta^{(t)})
  =\sum_z w_t(z)\log p_\theta(z\mid y).
\]
它不是前面的分布熵 $H(w_t)$；两者在当前点满足 $H(\theta^{(t)}\mid\theta^{(t)})=-H(w_t)$。由联合概率与条件概率的分解，
\[
  \ell(\theta)=Q(\theta\mid\theta^{(t)})-H(\theta\mid\theta^{(t)}),
\]
所以在各项有限时得到增量恒等式
\begin{equation}
  \begin{aligned}
    \ell(\theta^{(t+1)})-\ell(\theta^{(t)})
    &=Q(\theta^{(t+1)}\mid\theta^{(t)})-Q(\theta^{(t)}\mid\theta^{(t)})\\
    &\quad+D_{\mathrm{KL}}\!\left(w_t\Vert p_{\theta^{(t+1)}}(\,\cdot\mid y)\right).
  \end{aligned}
  \label{eq:lihang-em-increment-identity}
\end{equation}
在提高 $Q$ 的更新中，右边两项均非负；似然保持不变当且仅当 $Q$ 没有提高且新旧隐变量后验相同。因此 EM 的普遍保证是\emph{单调不减}，并不是每一步都严格增大，也不是一旦不变就已经得到唯一的极大值。

\subsection{目标值收敛、参数收敛与驻点}
\label{subsec:lihang-em-stationarity}

\begin{proposition}[目标值的收敛]
\label{prop:lihang-em-value-convergence}
设迭代从有限的 $\ell(\theta^{(0)})$ 出发，且存在有限常数 $M$ 使所有迭代点都满足 $\ell(\theta^{(t)})\leq M$。则 EM 产生的对数似然序列收敛到某个有限值 $\ell^*$。
\end{proposition}

\begin{proof}
由定理~\ref{thm:lihang-em-monotonicity}，该序列单调不减；再由给定上界，实数单调有界定理保证其收敛。
\end{proof}

这个命题只涉及函数值。相同似然可能对应不同参数，参数序列也可能没有唯一极限；例~\ref{ex:lihang-em-three-coins} 中的不可辨识性已经说明不能从最优似然唯一反推参数。对于连续分布，密度值没有概率质量不超过 $1$ 的限制，观测似然也未必有上界；高斯混合模型中就会出现这一问题。

教材关于参数极限为稳定点的结论需要额外的正则条件。下面给出一组便于核查的充分条件，只证明\emph{已经收敛的参数序列}具有什么性质。

\begin{proposition}[参数极限为驻点的一组充分条件]
\label{prop:lihang-em-stationary-limit}
设 $\Theta\subseteq\mathbb R^d$ 为开集，隐状态有限，$p_\theta(y,z)>0$ 且关于 $\theta$ 连续可微，$\nabla_\theta Q(\theta\mid\eta)$ 关于 $(\theta,\eta)$ 连续。若精确 EM 序列满足 $\theta^{(t)}\to\theta^*\in\Theta$，且各 M 步在域内取得最大值，则 $\nabla\ell(\theta^*)=0$。
\end{proposition}

\begin{proof}
域内的 M 步最大点满足 $\nabla_\theta Q(\theta^{(t+1)}\mid\theta^{(t)})=0$。令 $t\to\infty$，由梯度连续性得到 $\nabla_\theta Q(\theta^*\mid\theta^*)=0$。另一方面，有限求和允许逐项微分，故
\begin{equation}
  \begin{aligned}
    \nabla\ell(\theta)
    &=\frac{\sum_z\nabla p_\theta(y,z)}{p_\theta(y)}\\
    &=\sum_z p_\theta(z\mid y)\nabla\log p_\theta(y,z)
    =\left.\nabla_\vartheta Q(\vartheta\mid\theta)\right|_{\vartheta=\theta}.
  \end{aligned}
  \label{eq:lihang-em-score-identity}
\end{equation}
取 $\theta=\theta^*$，即得结论。
\end{proof}

这里的驻点（stationary point）仅表示梯度为零，可能是局部极大点、鞍点或退化的平坦点，不能直接理解为全局最优点。若参数受约束，应在可行方向上陈述相应的一阶必要条件；若隐变量连续，还必须增加使式~\eqref{eq:lihang-em-score-identity} 中微分与积分交换成立的条件。实际计算可采用多个可行初值，比较同一目标下得到的有限解，并检查约束、退化情形与停止状态；这些操作提高找到较好解的机会，但不构成一般的全局最优保证。

\section{EM 算法在高斯混合模型学习中的应用}
\label{sec:lihang-em-gmm}

\subsection{高斯混合模型}
\label{subsec:lihang-em-gmm-model}

\begin{definition}[一维高斯混合模型]
\label{def:lihang-em-gmm}
给定正整数 $K$，\emph{高斯混合模型}（Gaussian mixture model，GMM）的密度为
\begin{equation}
  p_\theta(y)=\sum_{k=1}^K\alpha_k\phi(y;\mu_k,\sigma_k^2),\qquad
  \phi(y;\mu_k,\sigma_k^2)
  =\frac1{\sqrt{2\pi\sigma_k^2}}
   \exp\left(-\frac{(y-\mu_k)^2}{2\sigma_k^2}\right),
  \label{eq:lihang-em-gmm-density}
\end{equation}
其中 $\alpha_k\geq0$、$\sum_{k=1}^K\alpha_k=1$、$\mu_k\in\mathbb R$、$\sigma_k^2>0$。参数 $\theta$ 包含各混合系数、均值与方差；成分数 $K$ 在本章的参数估计中预先给定。
\end{definition}

模型可用两步生成过程解释：先以概率 $\alpha_k$ 选择成分 $Z_i=k$，再按该成分的高斯分布生成 $Y_i$。假设 $n$ 个样本独立同分布，则
\[
  p_\theta(y_i,Z_i=k)=\alpha_k\phi(y_i;\mu_k,\sigma_k^2),\qquad
  \ell(\theta)=\sum_{i=1}^n\log\sum_{k=1}^K
  \alpha_k\phi(y_i;\mu_k,\sigma_k^2).
\]
若直接对观测对数似然求导，各参数会被成分求和耦合。引入成分指示随机变量 $W_{ik}=\mathbf1_{\{Z_i=k\}}$，其取值记为 $w_{ik}$；每个样本满足 $\sum_{k=1}^K w_{ik}=1$。完全数据的对数似然因此变为
\begin{equation}
  \ell_{\mathrm c}(\theta;y,w)
  =\sum_{i=1}^n\sum_{k=1}^K w_{ik}
  \left[\log\alpha_k-\frac12\log(2\pi\sigma_k^2)
  -\frac{(y_i-\mu_k)^2}{2\sigma_k^2}\right].
  \label{eq:lihang-em-gmm-complete-loglik}
\end{equation}
以下推导在混合系数为正、各方差为正的迭代点进行，因而观测密度及后验分母均为正。

\subsection{高斯混合模型参数估计的 EM 算法}
\label{subsec:lihang-em-gmm-updates}

E 步计算每个成分对每个观测的\emph{响应度}（responsibility）：
\begin{equation}
  \gamma_{ik}^{(t)}
  =\mathbb E_{\theta^{(t)}}[W_{ik}\mid Y=y]
  =\frac{\alpha_k^{(t)}\phi(y_i;\mu_k^{(t)},(\sigma_k^2)^{(t)})}
  {\sum_{h=1}^K\alpha_h^{(t)}\phi(y_i;\mu_h^{(t)},(\sigma_h^2)^{(t)})},\qquad
  \sum_{k=1}^K\gamma_{ik}^{(t)}=1.
  \label{eq:lihang-em-gmm-responsibility}
\end{equation}
这里 $\gamma_{ik}^{(t)}$ 是确定的条件概率，与尚未观测到的随机指示量 $W_{ik}$ 不同。记本轮的成分有效样本数为 $N_k^{(t)}=\sum_{i=1}^n\gamma_{ik}^{(t)}$，有 $\sum_kN_k^{(t)}=n$；它通常不是整数。将式~\eqref{eq:lihang-em-gmm-complete-loglik} 中的指示量换成其条件期望，得到
\begin{equation}
  Q(\theta\mid\theta^{(t)})
  =\sum_{k=1}^K\left[
  N_k^{(t)}\log\alpha_k-\frac{N_k^{(t)}}2\log(2\pi\sigma_k^2)
  -\frac1{2\sigma_k^2}\sum_{i=1}^n\gamma_{ik}^{(t)}(y_i-\mu_k)^2\right].
  \label{eq:lihang-em-gmm-q}
\end{equation}

M 步首先对混合系数进行约束极大化。对 $\sum_kN_k^{(t)}\log\alpha_k$ 引入约束 $\sum_k\alpha_k=1$ 的拉格朗日乘子，使用
\[
  \mathcal L(\alpha,\lambda)
  =\sum_kN_k^{(t)}\log\alpha_k+\lambda\left(1-\sum_k\alpha_k\right),
\]
一阶条件给出 $N_k^{(t)}/\alpha_k=\lambda$。求和得 $\lambda=n$，从而 $\alpha_k^{(t+1)}=N_k^{(t)}/n$。当所有 $N_k^{(t)}>0$ 时，目标关于混合系数严格凹，所以这是单纯形内的唯一最大点。

其次，对第 $k$ 个成分记 $S_k(\mu)=\sum_i\gamma_{ik}^{(t)}(y_i-\mu)^2$。固定方差时，极大化 $Q$ 等价于极小化 $S_k$，求导得到新均值
\[
  \overline y_k^{(t)}=\frac{\sum_i\gamma_{ik}^{(t)}y_i}{N_k^{(t)}}.
\]
恒等式 $S_k(\mu)=S_k(\overline y_k^{(t)})+N_k^{(t)}(\mu-\overline y_k^{(t)})^2$ 说明，对任意正方差，这个均值都使对应成分的 $Q$ 最大。再令 $v=\sigma_k^2$，在这个均值处有
\[
  \frac{\partial Q}{\partial v}
  =-\frac{N_k^{(t)}}{2v}+\frac{S_k(\overline y_k^{(t)})}{2v^2}.
\]
若 $S_k(\overline y_k^{(t)})>0$，导数在 $v=S_k/N_k^{(t)}$ 两侧由正变负，该点就是关于 $v>0$ 的唯一最大点。合并以上结果，得到
\begin{equation}
  \begin{aligned}
    \alpha_k^{(t+1)}&=\frac{N_k^{(t)}}n,\\
    \mu_k^{(t+1)}&=\frac{\sum_{i=1}^n\gamma_{ik}^{(t)}y_i}{N_k^{(t)}},\\
    (\sigma_k^2)^{(t+1)}
    &=\frac{\sum_{i=1}^n\gamma_{ik}^{(t)}(y_i-\mu_k^{(t+1)})^2}{N_k^{(t)}},
    \qquad k=1,\ldots,K.
  \end{aligned}
  \label{eq:lihang-em-gmm-m-step}
\end{equation}
\emph[caution]{新方差必须围绕本轮更新后的均值计算}。所有响应度均来自旧参数，但在同一次 M 步中，方差公式使用 $\mu_k^{(t+1)}$。分母是有效样本数 $N_k^{(t)}$，因为这里求极大似然估计，不能换成无偏样本方差中的自由度修正。若 $N_k^{(t)}=0$，该成分的均值和方差在本轮不可由 $Q$ 确定；若 $S_k=0$，正方差域内的最大值不存在，不能把方差零当作正常高斯分布。

具体计算时，输入样本、给定的 $K$、正混合系数与正方差的初值、容差和最大轮数；交替使用式~\eqref{eq:lihang-em-gmm-responsibility} 与式~\eqref{eq:lihang-em-gmm-m-step}，每轮计算观测对数似然并检查约束及停止条件，最后输出参数、似然与停止状态。一维情形下，一轮计算全部响应度和加权统计量需要 $O(nK)$ 次标量运算；若存储整个响应度矩阵，需要 $O(nK)$ 存储空间。密度很小时，可在对数域计算 $b_{ik}=\log\alpha_k^{(t)}+\log\phi(y_i;\mu_k^{(t)},(\sigma_k^2)^{(t)})$，令 $m_i=\max_k b_{ik}$，使用
\[
  \log\sum_k e^{b_{ik}}=m_i+\log\sum_k e^{b_{ik}-m_i}
\]
归一化响应度并求似然，以减轻直接计算很小密度造成的数值下溢。

\begin{example}
\label{ex:lihang-em-gmm-one-step}
对样本 $y=(-1,0,1)$ 拟合两个高斯成分。取初值 $\alpha_1^{(0)}=\alpha_2^{(0)}=1/2$、$\mu_1^{(0)}=-1$、$\mu_2^{(0)}=1$、两个方差都为 $1$，计算一轮 EM 更新。
\end{example}

\begin{solve}
两个初始密度的比值为 $\phi(y;-1,1)/\phi(y;1,1)=e^{-2y}$，所以第一成分的响应度依次为
\[
  \gamma_{11}^{(0)}=\frac{e^2}{1+e^2},\qquad
  \gamma_{21}^{(0)}=\frac12,\qquad
  \gamma_{31}^{(0)}=\frac1{1+e^2}.
\]
两个有效样本数均为 $3/2$，故混合系数不变。令 $a=(e^2-1)/(e^2+1)$，更新后的均值和方差为
\[
  \mu_1^{(1)}=-\frac{2a}{3}\approx-0.507729,\qquad
  \mu_2^{(1)}=\frac{2a}{3}\approx0.507729,
\]
\[
  (\sigma_1^2)^{(1)}=(\sigma_2^2)^{(1)}
  =\frac23-\frac{4a^2}{9}\approx0.408877.
\]
方差计算使用加权二阶矩减去新均值平方。该轮中间样本同时参与两个成分的估计，具体体现了响应度作为软分配权重的含义。
\end{solve}

\begin{notealert}[type=caution,title={高斯混合模型的似然可能无上界}]
当 $K\geq2$ 且各成分方差可任意趋近于零时，可以固定一个正的混合系数，将对应均值放在某个观测 $y_i$ 上，并令其标准差 $\sigma\to0$。该点的高斯密度为 $1/(\sqrt{2\pi}\sigma)\to\infty$，而另一个保持固定的正权重高斯成分可使其他观测的混合密度保持正的下界。因此总对数似然可趋于 $+\infty$，命题~\ref{prop:lihang-em-value-convergence} 的上界假设在无约束 GMM 中并不自动成立。

若模型明确约束 $\sigma_k^2\geq v_{\min}>0$，其中 $v_{\min}$ 是事先给定的方差下界，则对应的约束 M 步取 $\max\{v_{\min},S_k(\mu_k^{(t+1)})/N_k^{(t)}\}$，仍可按同一约束目标证明单调性；采用先验则对应另一种 MAP 目标。随意重置退化成分的参数未必提高原来的 $Q$，需要重新核查。即使排除了方差退化，交换成分编号仍不改变混合密度，因此类别编号本身也不具有唯一的统计含义。
\end{notealert}

\section{EM 算法的推广}
\label{sec:lihang-em-generalizations}

\subsection{$F$ 函数的极大--极大算法}
\label{subsec:lihang-em-f-function}

Jensen 下界中的后验权重也可以视为待优化的分布变量，从而将 EM 写成两个变量组之间的交替极大化。以下性质仍先在有限隐状态的概率单纯形上陈述；涉及局部极值的连续性结论不直接从有限情形推广到任意分布空间。

\begin{definition}[$F$ 函数]
\label{def:lihang-em-f-function}
对隐状态上的任意概率分布 $r=(r(z))$，定义其熵 $H(r)=-\sum_zr(z)\log r(z)$，并令
\begin{equation}
  F(r,\theta)=\mathbb E_r[\log p_\theta(y,Z)]+H(r)
  =\sum_zr(z)\log\frac{p_\theta(y,z)}{r(z)}.
  \label{eq:lihang-em-f-definition}
\end{equation}
在联合概率为正的有限状态情形，约定 $r(z)=0$ 的项为零。$r$ 是辅助分布，$\theta$ 是生成模型参数，二者承担不同的角色。
\end{definition}

\begin{lemma}[$F$ 的最大值与取等条件]
\label{lem:lihang-em-f-gap}
设 $0<p_\theta(y)<\infty$ 且联合概率在有限隐状态上为正，则
\begin{equation}
  \ell(\theta)-F(r,\theta)
  =D_{\mathrm{KL}}\!\left(r\Vert p_\theta(\,\cdot\mid y)\right)\geq0.
  \label{eq:lihang-em-f-gap}
\end{equation}
因此固定 $\theta$ 时，唯一使 $F$ 达到最大值的分布为 $r_\theta(z)=p_\theta(z\mid y)$，且
\begin{equation}
  F(r_\theta,\theta)=\ell(\theta)=\max_rF(r,\theta).
  \label{eq:lihang-em-f-profile}
\end{equation}
若 $p_\theta(y,z)$ 关于 $\theta$ 连续且分母保持为正，则 $r_\theta$ 关于 $\theta$ 连续。
\end{lemma}

\begin{proof}
代入 $p_\theta(y,z)=p_\theta(z\mid y)p_\theta(y)$，并利用 $\sum_zr(z)=1$，得到
\[
  F(r,\theta)=\ell(\theta)+\sum_zr(z)\log\frac{p_\theta(z\mid y)}{r(z)}.
\]
由 KL 散度的非负性及取等条件即得两式。后验分布每个坐标都是连续联合概率除以有限个连续联合概率之和，分母为正，因此连续。
\end{proof}

引理~\ref{lem:lihang-em-f-gap} 同时给出了教材两个引理的结论。也可用拉格朗日乘子核查最大分布：对 $F(r,\theta)+\lambda(1-\sum_zr(z))$ 关于正的 $r(z)$ 求导，得
\[
  \log p_\theta(y,z)-\log r(z)-1-\lambda=0.
\]
因而 $r(z)$ 与 $p_\theta(y,z)$ 成比例，归一化即得后验。熵项使 $F$ 关于 $r$ 严格凹，保证该解是唯一的最大分布，而不仅是一个满足一阶条件的候选点。

\begin{proposition}[联合极大点与似然极大点]
\label{prop:lihang-em-f-extrema}
在引理~\ref{lem:lihang-em-f-gap} 的有限状态、正联合概率及后验连续条件下，若 $(r^*,\theta^*)$ 是 $F$ 的联合局部极大点，则 $r^*=r_{\theta^*}$，且 $\theta^*$ 是 $\ell$ 的局部极大点。反过来，若 $\theta^*$ 是 $\ell$ 的局部极大点，则 $(r_{\theta^*},\theta^*)$ 是 $F$ 的联合局部极大点。把局部极大换成全局最大，两方向仍成立。
\end{proposition}

\begin{proof}
固定 $\theta^*$，$F$ 关于 $r$ 为严格凹函数。若 $r^*$ 不是其唯一最大点，沿着连接 $r^*$ 与 $r_{\theta^*}$ 的线段就会有任意近的点提高 $F$，与联合局部极大矛盾，所以 $r^*=r_{\theta^*}$。若存在趋向 $\theta^*$ 的参数序列使 $\ell$ 更大，由后验连续性，相应的 $(r_\theta,\theta)$ 趋向 $(r^*,\theta^*)$，但 $F(r_\theta,\theta)=\ell(\theta)>\ell(\theta^*)$，仍矛盾。反方向由 $F(r,\theta)\leq\ell(\theta)\leq\ell(\theta^*)=F(r_{\theta^*},\theta^*)$ 在相应邻域内成立得到。全局情形直接将邻域换成整个可行域。
\end{proof}

EM 的一次迭代恰好可以写为
\begin{equation}
  \begin{aligned}
    r^{(t+1)}&=p_{\theta^{(t)}}(\,\cdot\mid y)
      \in\operatorname*{arg\,max}_rF(r,\theta^{(t)}),\\
    \theta^{(t+1)}&\in\operatorname*{arg\,max}_{\theta\in\Theta}F(r^{(t+1)},\theta)
      =\operatorname*{arg\,max}_{\theta\in\Theta}Q(\theta\mid\theta^{(t)}).
  \end{aligned}
  \label{eq:lihang-em-max-max}
\end{equation}
第二行成立是因为固定 $r^{(t+1)}$ 后，$F=Q+H(r^{(t+1)})$ 中的熵不随 $\theta$ 变化。这就是教材的\emph{极大--极大算法}（maximization--maximization algorithm），亦即其 GEM 算法 1 所写的精确两步。M 步完成后，$r^{(t+1)}$ 仍是旧参数的后验，要到下一次 E 步才与新参数匹配。命题~\ref{prop:lihang-em-f-extrema} 谈的是两个变量同时变化下的联合极大点；仅仅分别对每组变量无法继续提高，并不自动满足该命题的联合局部极大假设。

\subsection{GEM 算法}
\label{subsec:lihang-em-gem}

\emph{广义期望极大算法}（generalized expectation maximization algorithm，GEM）保留 E 步的当前后验，但放宽精确 M 步：只要求新参数在可行域内满足
\begin{equation}
  Q(\theta^{(t+1)}\mid\theta^{(t)})
  \geq Q(\theta^{(t)}\mid\theta^{(t)}).
  \label{eq:lihang-em-gem-condition}
\end{equation}
未停止时通常寻求严格提高。由定理~\ref{thm:lihang-em-monotonicity}，这种更新仍使观测似然不减；但放宽条件也可能容许没有实际进展的更新，所以仅凭该不等式不能证明参数收敛或其极限为驻点。教材的 GEM 算法 2 就是每轮计算 $Q$ 后，寻找满足上述提高条件的参数，继而检查停止准则。

例如在无约束开域中，若本轮 $Q$ 在当前点可微且 $g_t=\nabla_\theta Q(\theta^{(t)}\mid\theta^{(t)})\ne0$，则沿该方向的一阶展开为
\[
  Q(\theta^{(t)}+\eta g_t\mid\theta^{(t)})
  =Q(\theta^{(t)}\mid\theta^{(t)})+\eta\|g_t\|^2+o(\eta).
\]
因而充分小的正步长会提高 $Q$；可通过逐步缩短步长并实际检查目标增量来选取更新。存在参数约束时，还需使整个更新可行，不能直接沿任意梯度方向越出概率单纯形或正方差域。

教材的 GEM 算法 3 将 M 步分成若干次\emph{条件极大化}（conditional maximization）。设参数分为 $d$ 个坐标或参数块，令 $\theta^{(t,0)}=\theta^{(t)}$。在第 $j$ 次内层更新中，只改变第 $j$ 块，求
\begin{equation}
  \theta_j^{(t,j)}\in\operatorname*{arg\,max}_{u:
    (\theta_1^{(t,j-1)},\ldots,u,\ldots,\theta_d^{(t,j-1)})\in\Theta}
  Q(\theta_1^{(t,j-1)},\ldots,u,\ldots,\theta_d^{(t,j-1)}\mid\theta^{(t)}),
  \label{eq:lihang-em-conditional-maximization}
\end{equation}
其余块保持 $\theta_h^{(t,j)}=\theta_h^{(t,j-1)}$，$h\ne j$。依次进行 $j=1,\ldots,d$，再置 $\theta^{(t+1)}=\theta^{(t,d)}$。每个条件可行域都包含更新前的该块值，因此每次内层更新不降低 $Q$，串联即得式~\eqref{eq:lihang-em-gem-condition}。\emph[caution]{整轮条件极大化固定同一个旧后验}；已经更新的参数块用于后续条件优化，但不能把后验权重也无说明地逐块换掉。各条件最大值必须存在，带有耦合约束的坐标若无法单独改变，需选择合适的参数块。

上述放宽直接作用于 M 步。若同时限制辅助分布的形式、令 E 步不能取得真实后验，则通常只能从交替提高 $F$ 得到辅助下界不减；原似然的单调性还需要额外论证。原因是此时 $F$ 与 $\ell$ 之间的 KL 差距不一定在每轮起点为零。这也再次说明，EM 的似然保证依赖当前后验给出的接触条件。

\section*{本章回顾与教材习题}
\addcontentsline{toc}{section}{本章回顾与教材习题}

EM 将观测似然的优化分解为后验计算和完全数据条件期望的优化；Jensen 下界、增量恒等式与 $F$ 函数是同一机制的三种表达。具体模型中的推导应首先说明联合分布与隐变量，其次计算后验权重，最后验证带约束的 M 步是否真正存在并提高目标。单调性、目标值收敛、参数收敛及全局最优是不同层次的结论。教材习题 9.2 要求的 $F(r_\theta,\theta)=\ell(\theta)$ 已在引理~\ref{lem:lihang-em-f-gap} 中证明，以下补充其余三题的计算与构造。

\begin{example}
\label{ex:lihang-em-coin-exercise}
教材习题 9.1。沿用例~\ref{ex:lihang-em-three-coins} 的观测数据，以 $(\pi^{(0)},p^{(0)},q^{(0)})=(0.46,0.55,0.67)$ 为初值求参数估计，并解释其含义。
\end{example}

\begin{solve}
对于正面和反面观测，分别使用
\[
  r_1=\frac{0.46\times0.55}{0.46\times0.55+0.54\times0.67},\qquad
  r_0=\frac{0.46\times0.45}{0.46\times0.45+0.54\times0.33}.
\]
代入式~\eqref{eq:lihang-em-coin-update}，第一轮得到
\[
  (\pi^{(1)},p^{(1)},q^{(1)})
  \approx(0.461863,0.534595,0.656135).
\]
更一般地，只要更新分母为正，令 $\overline y=n^{-1}\sum_i y_i$，同一公式直接给出
\[
  \pi^{(t+1)}p^{(t+1)}
  +(1-\pi^{(t+1)})q^{(t+1)}
  =\frac{\sum_i r_i^{(t)}y_i+\sum_i(1-r_i^{(t)})y_i}{n}
  =\overline y.
\]
本题第一轮已使混合正面概率等于 $0.6$，故到达极大似然解集。由于这类观测的响应度只取 $r_1,r_0$ 两个值，下一轮正面响应度为 $\pi^{(1)}p^{(1)}/\overline y=r_1$，反面响应度同样为 $r_0$，所以参数随后保持不变。不同初值所得参数可以不同，却有相同的最优观测似然。
\end{solve}

\begin{example}
\label{ex:lihang-em-gmm-exercise}
教材习题 9.3。对数据
\[
  -67,-48,6,8,14,16,23,24,28,29,41,49,56,60,75
\]
拟合两个成分的一维高斯混合模型，给出一次具有明确初值和停止条件的数值估计。
\end{example}

\begin{solve}
取初始混合系数 $(0.5,0.5)$、均值 $(-50,30)$、方差 $(400,400)$，按式~\eqref{eq:lihang-em-gmm-responsibility} 和式~\eqref{eq:lihang-em-gmm-m-step} 进行精确形式的 EM 更新。在对数域求响应度及观测似然，设最大迭代次数为 $1000$，连续两轮观测对数似然绝对差小于 $10^{-10}$ 时停止。本次计算在第 $5$ 轮满足该条件，结果约为
\[
  \begin{aligned}
    (\widehat\alpha_1,\widehat\alpha_2)&=(0.133172,0.866828),\\
    (\widehat\mu_1,\widehat\mu_2)&=(-57.511077,32.984887),\\
    (\widehat\sigma_1^2,\widehat\sigma_2^2)&=(90.249878,429.458341),
  \end{aligned}
\]
相应观测对数似然约为 $-71.063362$，计算中各轮目标值不减。这里报告的是该初值和停止规则下的有限数值迭代结果，其全局最优性未经证明；无约束模型的似然退化问题仍然存在。交换两个成分的编号会得到同一个混合密度，比较不同运行时应先对齐编号。
\end{solve}

\begin{example}
\label{ex:lihang-em-naive-bayes}
教材习题 9.4。设样本 $x_1,\ldots,x_n$ 没有类别标记，各 $x_i=(x_{i1},\ldots,x_{im})$ 由离散特征组成，第 $j$ 个特征的有限取值集合为 $\mathcal A_j$。写出给定类别数 $K$ 时，朴素贝叶斯模型无监督参数估计的 EM 算法。
\end{example}

\begin{solve}
令潜在类别为 $C_i\in\{1,\ldots,K\}$，类别先验为 $\pi_k$，特征条件概率为 $\beta_{jka}=P(X_{ij}=a\mid C_i=k)$。模型假设各样本独立同分布，且给定类别后各特征条件独立，于是
\[
  p_\theta(x_i,C_i=k)=\pi_k\prod_{j=1}^m\beta_{jkx_{ij}},\qquad
  \sum_k\pi_k=1,\qquad \sum_{a\in\mathcal A_j}\beta_{jka}=1.
\]
输入观测样本、类别数、各特征取值集合及正的初始概率，E 步计算
\[
  \tau_{ik}^{(t)}
  =\frac{\pi_k^{(t)}\prod_{j=1}^m\beta_{jkx_{ij}}^{(t)}}
  {\sum_{h=1}^K\pi_h^{(t)}\prod_{j=1}^m\beta_{jhx_{ij}}^{(t)}}.
\]
完全数据对数似然中的类别指示量取条件期望后，得到
\[
  Q(\theta\mid\theta^{(t)})
  =\sum_{i=1}^n\sum_{k=1}^K\tau_{ik}^{(t)}
  \left[\log\pi_k+\sum_{j=1}^m\log\beta_{jkx_{ij}}\right].
\]
对每组概率施加归一化约束，M 步即为后验加权频率：
\[
  N_k^{(t)}=\sum_i\tau_{ik}^{(t)},\qquad
  \pi_k^{(t+1)}=\frac{N_k^{(t)}}n,\qquad
  \beta_{jka}^{(t+1)}
  =\frac{\sum_i\tau_{ik}^{(t)}\mathbf1_{\{x_{ij}=a\}}}{N_k^{(t)}}.
\]
在 $N_k^{(t)}>0$ 时，这些式子分别是相应多项分布似然的最大点；若有效样本数为零，该类别的条件参数不由本轮目标确定。交替 E、M 两步，按观测对数似然容差或最大轮数停止，输出类别先验与各特征条件概率。应避免所有类别使用完全相同的初始参数，否则对称性可能一直保持；没有真实类别标记时，潜在类别的编号不能直接解释为预先指定的语义类别。
\end{solve}

% 教材来源：机器学习/textbooks/机器学习（分章原文）/机器学习-10.pdf。
% 对应教材第 10 章，原书第 193--214 页；本文件采用自动章节编号。




